% !TeX program = lualatex
%==============================================================================
%  Leant — a GHCi-style REPL for Lean 4
%  An overview and tutorial, with a detailed exposition of :synth.
%
%  Leant is an experimental project in active development; this document
%  describes the state of the tool at the date on the title page.
%
%  Build:  lualatex Leant.tex   (twice, for the table of contents)
%
%  The visual design (palette, fonts, code boxes) follows the typeset
%  edition of "Counterexamples in Type Systems".
%==============================================================================
\documentclass[11pt,oneside]{article}

\usepackage[letterpaper,
            textwidth=5.65in, textheight=8.5in,
            centering, headsep=16pt, footskip=28pt]{geometry}

\usepackage{amsmath}
\usepackage{fontspec}
\usepackage{unicode-math}
\usepackage{microtype}
\usepackage{xcolor}
\usepackage{fvextra}
\usepackage[most]{tcolorbox}
\usepackage{enumitem}
\usepackage{titlesec}
\usepackage{fancyhdr}
\usepackage{needspace}
\usepackage[hang,flushmargin,bottom]{footmisc}
\usepackage{xurl}
\usepackage{hyperref}
\usepackage{bookmark}

%------------------------------------------------------------------ fonts ----
\setmainfont{LibertinusSerif}[
  Extension    = .otf,
  UprightFont  = *-Regular, ItalicFont = *-Italic,
  BoldFont     = *-Bold,    BoldItalicFont = *-BoldItalic,
  Ligatures    = TeX ]
\setsansfont{LibertinusSans}[
  Extension    = .otf,
  UprightFont  = *-Regular, ItalicFont = *-Italic, BoldFont = *-Bold,
  Ligatures    = TeX ]
\setmonofont{DejaVuSansMono}[
  Extension    = .ttf,
  UprightFont  = *, ItalicFont = *-Oblique,
  BoldFont     = *-Bold, BoldItalicFont = *-BoldOblique,
  Scale        = 0.83 ]
\setmathfont{LibertinusMath-Regular.otf}

%----------------------------------------------------------------- colours ---
%  A calm, paper-like palette: sage greens, warm off-whites, near-black.
\definecolor{ink}      {HTML}{0B0B0B}  % body text
\definecolor{moss}     {HTML}{3E5C43}  % links, rules, section numbers
\definecolor{sage}     {HTML}{D7E8CD}  % code label ground
\definecolor{sagepale} {HTML}{F8FBF1}  % code ground
\definecolor{sageline} {HTML}{DFE4D7}  % hairlines
\definecolor{parchment}{HTML}{F4F3EB}  % warm ground
\definecolor{stone}    {HTML}{ADADAC}  % dimmed text
\definecolor{lichen}   {HTML}{77836A}  % tag line
\definecolor{cmtcol}   {HTML}{7A8A76}  % code comments
\definecolor{kwcol}    {HTML}{33523B}  % code keywords
\definecolor{strcol}   {HTML}{8A5A44}  % code strings
\color{ink}

%---------------------------------------------------------------- headings ---
\titleformat{\section}
  {\normalfont\LARGE\bfseries\color{ink}}
  {\color{moss}\thesection}{0.7em}{}
  [\vspace{2pt}{\color{sageline}\titlerule[0.8pt]}]
\titlespacing*{\section}{0pt}{0pt}{1.1\baselineskip}

\titleformat{\subsection}{\normalfont\large\bfseries\color{ink}}
  {\color{moss}\thesubsection}{0.6em}{}
\titlespacing*{\subsection}{0pt}{1.4\baselineskip}{0.45\baselineskip}

\setcounter{secnumdepth}{2}
\setcounter{tocdepth}{2}

%----------------------------------------------------------------- running ---
\pagestyle{fancy}
\fancyhf{}
\renewcommand{\headrulewidth}{0pt}
\fancyhead[L]{\footnotesize\itshape\color{stone}\nouppercase{\leftmark}}
\fancyhead[R]{\footnotesize\color{stone}\thepage}
\fancypagestyle{plain}{\fancyhf{}%
  \fancyhead[R]{\footnotesize\color{stone}\thepage}}
\renewcommand{\sectionmark}[1]{\markboth{#1}{}}

%------------------------------------------------------------------- lists ---
\setlist[itemize]{leftmargin=1.35em, itemsep=0.45\baselineskip,
                  parsep=0.35\baselineskip, topsep=0.5\baselineskip,
                  label={\color{moss}\rule[0.28ex]{0.36em}{0.36em}}}

%------------------------------------------------------------------- misc ----
\newcommand{\code}[1]{\texttt{#1}}
\newcommand{\cmd}[1]{\texttt{\color{kwcol}#1}}

\widowpenalty=10000
\clubpenalty=10000
\raggedbottom
\setlength{\emergencystretch}{2em}

%---------------------------------------------------------------- code box ---
%  REPL transcripts are typeset by fvextra rather than listings: the
%  transcripts are full of non-ASCII Lean syntax, and listings under
%  LuaTeX reorders a multibyte character in front of whatever precedes
%  it.  The session environment writes its body verbatim to a scratch
%  file and reads it back inside the box.
\tcbset{cebox/.style={
  enhanced, breakable,
  colback=sagepale, colframe=sageline, boxrule=0.5pt, arc=2pt,
  left=7pt, right=7pt, top=5pt, bottom=5pt,
  toptitle=1pt, bottomtitle=1pt,
  coltitle=ink, colbacktitle=sage,
  fonttitle=\sffamily\scriptsize\bfseries,
  attach boxed title to top left={xshift=9pt, yshift=-\tcboxedtitleheight/2},
  boxed title style={boxrule=0pt, arc=1.5pt, left=5pt, right=5pt,
                     top=1.5pt, bottom=1.5pt},
  before skip=\medskipamount, after skip=\medskipamount }}

\newtcolorbox{sessionouter}{cebox, title={session}}

\newenvironment{session}
 {\VerbatimEnvironment\begin{VerbatimOut}{\jobname.session.tmp}}
 {\end{VerbatimOut}%
  \begin{sessionouter}%
  \VerbatimInput[fontsize=\footnotesize, breaklines=true,
    breaksymbolleft={\textcolor{stone}{\footnotesize$\hookrightarrow$}},
    formatcom=\color{ink}, tabsize=4]{\jobname.session.tmp}%
  \end{sessionouter}}

% the "experimental" notice
\newtcolorbox{notice}[1]{enhanced, breakable,
  colback=parchment, colframe=sageline, boxrule=0.5pt, arc=2pt,
  left=8pt, right=8pt, top=6pt, bottom=6pt,
  coltitle=ink, colbacktitle=sage,
  fonttitle=\sffamily\scriptsize\bfseries, title={#1},
  attach boxed title to top left={xshift=9pt, yshift=-\tcboxedtitleheight/2},
  boxed title style={boxrule=0pt, arc=1.5pt, left=5pt, right=5pt,
                     top=1.5pt, bottom=1.5pt},
  before skip=\medskipamount, after skip=\medskipamount }

%------------------------------------------------------------ command table --
\newenvironment{cmdtable}
  {\par\begingroup\small
   \setlength{\parindent}{0pt}\setlength{\parskip}{2.5pt}}
  {\endgroup\par\vspace{0.3\baselineskip}}
\newcommand{\cmdrow}[2]{%
  \makebox[11.5em][l]{\code{#1}}\parbox[t]{\dimexpr\linewidth-12em}{#2}\par}

%--------------------------------------------------------------- hyperlinks --
\hypersetup{
  colorlinks = true,
  linkcolor  = moss,
  urlcolor   = moss,
  citecolor  = moss,
  pdftitle   = {Leant — a GHCi-style REPL for Lean 4},
  pdfauthor  = {The ProveIt project},
  pdfsubject = {An interactive REPL for Lean 4, with verified term synthesis},
  pdfkeywords= {Lean 4, REPL, term synthesis, Djinn, LJT, Exference},
  bookmarksnumbered = true }
\urlstyle{same}

\begin{document}
%------------------------------------------------------------------ title ---
\begin{titlepage}
\thispagestyle{empty}
\vspace*{\stretch{2}}
\begin{center}
  {\color{sageline}\rule{0.55\textwidth}{0.8pt}}\\[1.7em]
  {\fontsize{34}{40}\selectfont Leant}\\[0.55em]
  {\Large a GHCi-style REPL for Lean 4}\\[1.5em]
  {\color{sageline}\rule{0.55\textwidth}{0.8pt}}\\[2.2em]
  {\large \emph{an overview and tutorial, with a detailed tour of \code{:synth}}}
\end{center}
\vspace{\stretch{3}}
\begin{center}
  \small\color{lichen}
  An \textbf{experimental} tool in active development --- interfaces,
  output, and capabilities change frequently.\\[0.6em]
  Part of the \emph{ProveIt} repository, \code{Tools/Leant}.\quad
  July 2026.
\end{center}
\vspace*{\stretch{1}}
\end{titlepage}

\clearpage
{\hypersetup{linkcolor=ink}%
 \microtypesetup{protrusion=false}%
 \tableofcontents
 \microtypesetup{protrusion=true}}
\clearpage

%==============================================================================
\section{Introduction}\label{sec:intro}

Leant is an interactive read--eval--print loop for Lean~4, in the
spirit of GHCi: you type expressions and they are evaluated, you type
declarations and they enter the session, and a family of colon-commands
gives you type queries, documentation, search, interactive proving, and
--- the subject of half of this document --- automatic \emph{term
synthesis} with machine-checked answers.

Lean~4 itself is normally driven from an editor: the language server
shows goals and diagnostics as you edit a file. Leant complements that
workflow with a conversational one. It is aimed at exploration ---
trying a lemma, poking at a definition, asking \emph{is this type even
inhabited?} --- where the unit of work is a line, not a file.

\begin{notice}{experimental}
Leant is an \textbf{experimental project in active development}. It is
a research vehicle, not a supported product: commands appear, change
shape, and disappear between commits; output formats are not stable;
and some features are deliberately bounded approximations. Every
transcript in this document was produced by the tool at the date on
the title page and may not match the tool you build tomorrow. When
this document and \code{:help} disagree, believe \code{:help}.
\end{notice}

Two implementations live side by side in the repository:
\code{Tools/LeantPy} (the Python original, whose backend management is
delegated to LeanInteract) and \code{Tools/Leant} (the Haskell port,
now the primary implementation, which speaks the backend protocol
directly). They share the command set and the same session semantics;
term synthesis (\S\ref{sec:synth}) is Haskell-only, because it embeds
the Djex synthesis library in-process. Everything below describes the
Haskell implementation.

Under the hood, Leant drives the community
\code{leanprover-community/repl} binary over a JSON protocol: every
input becomes a command against a persistent Lean environment, and
every claim Leant makes --- including every synthesized term --- is
checked by that backend before it is shown. Sessions survive backend
crashes (the session history replays automatically), can be saved and
restored (\code{:pickle}/\code{:unpickle}), and can be recorded as
transcripts.

%==============================================================================
\section{Getting started}\label{sec:start}

Build the binary once with \code{cabal build exe:leant} (GHC 9.12.4;
the vendored \code{lib/Djex} submodule must be initialized), then run
\code{leant.cmd} --- or the built executable directly. Leant finds the
Lean backend the same way LeanInteract's cache laid it out, or accepts
an explicit \code{-{}-repl-exe}/\code{LEANT\_BACKEND}. Without a Lake
project it runs against the plain Lean~4 standard library, which is
also the least memorable thing about it: no imports, subsecond
startup.

A session begins as a calculator and escalates quietly:

\begin{session}
λ> 2 + 2
4
λ> it * 10
40
λ> def double (n : Nat) : Nat := n + n
λ> double 21
42
λ> :type double
double : Nat → Nat
\end{session}

Expressions are tried with \code{\#eval} and fall back to
\code{\#check}; the last evaluated value is bound as \code{it},
GHCi-style. Declarations (\code{def}, \code{theorem},
\code{inductive}, \dots) extend the session environment, and any
\code{\#}-command (\code{\#eval}, \code{\#check}, \code{\#print
axioms}) passes straight through. Multi-line input opens automatically
when a line is syntactically incomplete --- type a \code{structure}
header and keep going --- and an empty line submits it.

%==============================================================================
\section{The commands}\label{sec:commands}

\subsection{Session and evaluation}
\begin{cmdtable}
\cmdrow{:help, :h, :?}{show the command summary}
\cmdrow{:quit, :q}{exit}
\cmdrow{:type EXPR, :t}{show the type of \code{EXPR} (\code{\#check})}
\cmdrow{:info NAME, :i}{show the definition of \code{NAME} (\code{\#print})}
\cmdrow{:load FILE, :l}{reset the session and load a \code{.lean} file}
\cmdrow{:reload, :r}{reload the last loaded file}
\cmdrow{:import MOD}{add an import (rebuilds the session)}
\cmdrow{:imports}{list active imports}
\cmdrow{:undo}{revert the last state-changing command}
\cmdrow{:reset}{clear all definitions (keeps imports)}
\cmdrow{:history}{show the session's state-changing commands}
\cmdrow{:set OPT VAL}{a Lean \code{set\_option} that persists in the session}
\end{cmdtable}

\subsection{Discovery}
\begin{cmdtable}
\cmdrow{:browse [NS]}{list declarations in a namespace, or the session's own}
\cmdrow{:browse!\ NS}{\dots including compiler-generated auxiliaries}
\cmdrow{:doc NAME}{show a docstring}
\cmdrow{:search TEXT}{search declaration names, case-insensitively}
\cmdrow{:search?\ TYPE}{proof search: what proves \code{TYPE}? (via \code{exact?})}
\end{cmdtable}

\subsection{Proving and synthesis}
\begin{cmdtable}
\cmdrow{:prove [PROP]}{prove \code{PROP} interactively, tactic by tactic
  (without an argument: resume the last \code{sorry})}
\cmdrow{:synth TYPE}{synthesize verified terms of \code{TYPE};
  candidates are bound as \code{it1} (= \code{it}), \code{it2}, \dots
  (\S\ref{sec:synth})}
\cmdrow{:set synth-engine E}{\code{djinn} (default) $\mid$
  \code{exference} $\mid$ \code{both}}
\cmdrow{:set synth-steps N}{Exference's step budget (default 4096)}
\cmdrow{:set synth-classical B}{offer classical candidates for
  constructively refuted goals (\code{on}$\mid$\code{off}, default on)}
\end{cmdtable}

\subsection{Persistence and utility}
\begin{cmdtable}
\cmdrow{:transcript [F|on|off]}{record a transcript of the session}
\cmdrow{:timestamps [on|off]}{timestamp transcript entries}
\cmdrow{:pickle FILE}{save the environment to \code{FILE} (\code{.olean})}
\cmdrow{:unpickle FILE}{restore an environment}
\cmdrow{:env}{show the backend environment id}
\cmdrow{:time}{toggle per-command timing}
\cmdrow{:!\ CMD}{run a shell command}
\end{cmdtable}

%==============================================================================
\section{Interactive proving}\label{sec:prove}

\code{:prove} turns the prompt into a tactic loop against the
backend's proof-state protocol. Every line is a tactic; \code{:undo}
pops one step; \code{:qed NAME} assembles the accumulated script into
a theorem and replays it into the session, so the result is usable
afterwards.

\begin{session}
λ> :prove ∀ p q : Prop, (p → q) → p → q ∧ p
⊢ ∀ (p q : Prop), (p → q) → p → q ∧ p
⊢> intro p q h hp
p q : Prop
h : p → q
hp : p
⊢ q ∧ p
⊢> exact ⟨h hp, hp⟩
All goals accomplished
⊢> :qed mp_and
saved: theorem mp_and : ∀ p q : Prop, (p → q) → p → q ∧ p
\end{session}

A \code{sorry} in an ordinary declaration prints its goal and offers
the same loop via bare \code{:prove}. Prove mode also cooperates with
\code{:synth}, which is where it gets interesting
(\S\ref{sec:synthprove}).

%==============================================================================
\section{Term synthesis: \code{:synth}}\label{sec:synth}

\code{:synth TYPE} answers the question \emph{``write me a term of
this type''} --- or, on the propositions-as-types side of the mirror,
\emph{``prove this''} --- and sometimes the stronger question
\emph{``show me that no such term exists.''} It is built on the
vendored Djex library, linked in-process. Djex began by merging two
classic Haskell program synthesizers behind one checked synthesis
foundation --- \textbf{Djinn}, a complete and terminating proof
search for intuitionistic propositional logic (Dyckhoff's LJT
calculus) that emits programs, and \textbf{Exference}, a ranked
heuristic search under explicit budgets --- and has since moved well
beyond both originals. It fixes a substantial number of their bugs
(notably verdict honesty: an exhausted search over an approximated
space is reported as \emph{no evidence}, never as a refutation);
it implements markedly better type-class support, with a unified,
validated constraint contract shared by both engines --- Exference
resolves givens, superclasses, and instances, Djinn checks contexts
and synthesizes dictionary-independent terms; and it adds a growing
level of support for rank-N and impredicative types: goal-side
quantifiers open through polarity-aware plan families, quantified
hypotheses are instantiated at a bounded set of sequent-supplied
types, and impredicative instantiation is admitted under a guard that
never invents a polytype the query did not supply. That scope is
still being improved; Leant consumes Djex as a library behind a
deliberately narrow engine boundary, so each improvement arrives by
version bump. The engines speak Haskell types; Leant translates Lean
goals into their fragment and their answers back into Lean.

Three design rules keep this honest:

\begin{itemize}
\item \textbf{The engine is never trusted.} Every candidate is
  re-elaborated by the Lean backend against the exact goal
  (\code{example :\ (T) :=\ term}) before you see it. A rendering bug
  or an engine bug costs a dropped candidate, never a wrong answer.
\item \textbf{Refusals are honest.} Goals outside the supported
  fragment are refused with a reason, not silently mangled.
\item \textbf{Negative verdicts are labeled by strength.}
  ``Provably uninhabited'' appears only when the translation was
  complete and lossless; anything weaker is reported as ``no term
  found within bounds,'' which claims nothing.
\end{itemize}

\subsection{Higher-order plumbing}\label{sec:plumbing}

The sweet spot is the structural fragment: arrows, products, sums,
$\bot$, $\top$, negation, \code{Iff}, and quantifiers over opaque
variables --- the ``plumbing'' terms one writes constantly.
Free capital identifiers are auto-bound, so quick queries stay quick:

\begin{session}
λ> :synth ((A → B → C) → (A → B) → A → C)
  it1  fun f g x => f x (g x)
λ> :synth (A × (B ⊕ C) → (A × B) ⊕ (A × C))
  it1  fun ⟨x, y⟩ => match y with | .inl z => .inl ⟨x, z⟩ | .inr w => .inr ⟨x, w⟩
λ> :synth (∀ p q : Prop, (p ↔ q) → ¬p → ¬q)
  it1  fun _ _ ⟨_, f⟩ k x => k (f x)
\end{session}

The first is the S combinator; the second distributes a product over a
sum; the third is contraposition across an \code{Iff}, with the
backward implication projected out by the pattern
\code{⟨\_, f⟩}. Binders are named by role --- functions \code{f g h},
values \code{x y z}, negations and continuations \code{k} --- which is
a small thing that makes large candidates legible.

Candidates are ranked smallest-first and \emph{bound into the session}
as \code{it1}, \code{it2}, \dots, with bare \code{it} the best one.
They are ordinary definitions; you can evaluate them immediately:

\begin{session}
λ> :synth (a → a → a)
  it1  fun _ x => x
  it2  fun x _ => x
λ> #eval it2 "left" "right"
"left"
\end{session}

Both inhabitants of $a \to a \to a$ that matter, and picking one is
just using its name. For \emph{programs} --- as opposed to
proof-irrelevant propositions --- candidate choice is the whole point.

\subsection{Impossibility, proved}\label{sec:negative}

When Djinn's complete search exhausts a fully translated goal, failure
is a theorem:

\begin{session}
λ> :synth (∀ a b : Type, a → b)
provably uninhabited — no closed term of this polymorphic type exists
\end{session}

No Lean tactic gives you this answer (a failing \code{exact?} proves
nothing). Note the careful wording: the verdict is about \emph{closed
terms of the polymorphic type}. It does not say any particular
instance $a \to b$ is empty --- choose $a = b$ and it isn't --- and
for propositions it is about \emph{constructive} provability, a
distinction \S\ref{sec:classical} exploits. When the translation had
to hide structure behind an opaque atom, the verdict honestly
downgrades:

\begin{session}
λ> :synth (∀ a : Type, Sigma (fun _ : Nat => a))
no term found within bounds (opaque atoms `Nat` hide structure the engine cannot analyze — not a refutation)
\end{session}

\subsection{Inductive types}\label{sec:inductives}

A non-recursive, non-indexed inductive or structure applied to all its
parameters is expanded into a generalized sum of products:
constructors become introduction rules, case analysis the elimination
rule. This covers the standard library's little workhorses ---

\begin{session}
λ> :synth Ordering
  it1  .lt
  it2  .eq
  it3  .gt
λ> :synth (∀ a : Type, Option (Option a) → Option a)
  it1  fun _ _ => .none
  it2  fun _ x => match x with | .none => (.none) | .some y => y
λ> :synth (∀ e a b : Type, Except e a → (a → b) → Except e b)
  it1  fun _ _ _ x f => match x with | .error y => .error y | .ok z => .ok (f z)
\end{session}

--- \code{Option.join} and \code{Except.map} synthesized rather than
remembered. It extends to \code{Type}-valued classes-as-data like
\code{Decidable}, where the eliminations write themselves:

\begin{session}
λ> :synth (∀ p : Prop, Decidable p → Decidable (¬ p))
  it1  fun _ x => match x with | .isFalse k => .isTrue k | .isTrue y => .isFalse (fun k1 => k1 y)
\end{session}

Session-declared types participate the moment you declare them, and
refutations over expanded inductives remain sound, because the engine
saw the complete constructor list:

\begin{session}
λ> structure Pair (A B : Type) where
…>   fst : A
…>   snd : B
…>
λ> :synth (∀ a b : Type, a → b → Pair a b)
  it1  fun _ _ x y => ⟨x, y⟩
λ> :synth (∀ a b : Type, Pair a b → Empty)
provably uninhabited — no closed term of this polymorphic type exists
\end{session}

Recursive types cannot be expanded --- elimination is where
undecidability lives --- but their \emph{constructors} are still sound
introduction rules, so building is possible even where analysis is
not:

\begin{session}
λ> :synth Nat
  it1  Nat.zero
  it2  Nat.succ Nat.zero
λ> :synth (∀ a : Type, a → List a)
  it1  fun _ _ => List.nil
\end{session}

\subsection{Classical candidates}\label{sec:classical}

Peirce's law has no constructive inhabitant, and \code{:synth} can
prove that. But rather than stopping at the refutation, it then asks
the classical question, and answers with a term whose spelling
announces exactly what was used:

\begin{session}
λ> :synth (∀ p q : Prop, ((p → q) → p) → p)
  it1  fun _ _ f => match Classical.em _ with | .inl x => x | .inr k => f (fun y => absurd y k)
λ> :synth (∀ p : Prop, ¬¬p → p)
  it1  fun _ k => match Classical.em _ with | .inl x => x | .inr k1 => absurd k1 k
\end{session}

These are the proofs a human would write: one case split on excluded
middle, one \code{absurd} where the contradiction lands. Even
statements that are classically true but intuitionistically
unprovable in a less famous way come out as case splits --- the
G\"odel--Dummett formula, for instance:

\begin{session}
λ> :synth (∀ p q : Prop, (p → q) ∨ (q → p))
  it1  fun _ _ => match Classical.em _ with | .inl x => (match Classical.em _ with | .inl _ => .inr (fun _ => x) | .inr k => .inr (fun y => absurd y k)) | .inr k1 => .inl (fun z => absurd z k1)
\end{session}

Two searches stand behind this. First, one excluded-middle assumption
per atomic subformula is handed to the intuitionistic engine ---
complete for the propositional fragment, since intuitionistic logic
plus atom instances of $p \vee \neg p$ \emph{is} classical
propositional logic --- and rendered as \code{Classical.em} case
splits. If that misses, the goal's double-negation translation runs
instead (complete by Glivenko's theorem) and the candidate is wrapped
in \code{Classical.byContradiction}. Both kinds are verified like any
other candidate. The whole behavior sits behind a switch:

\begin{session}
λ> :set synth-classical off
synth classical: off
λ> :synth (∀ p : Prop, p ∨ ¬ p)
provably uninhabited — no closed term of this polymorphic type exists
(constructively — a classical proof may still exist; this is not a disproof of the proposition)
\end{session}

\subsection{Dependent formulas as cargo}\label{sec:dependent}

Leant's translation deliberately does not analyze dependent types.
But it does not refuse them either: a dependent subformula becomes an
opaque \emph{atom}, compared up to $\alpha$-equivalence, which the
engine can carry from one side of a goal to the other. Transport
works; analysis does not:

\begin{session}
λ> opaque P : Nat → Prop
λ> opaque Q : Prop
λ> :synth ((∀ n : Nat, P n) ∧ Q → Q ∧ (∀ n : Nat, P n))
  it1  fun ⟨x, y⟩ => ⟨y, x⟩
\end{session}

The engine never looked inside $\forall n,\, P\,n$ --- it swapped a
sealed box. A goal that would require opening the box (an induction, a
rewrite, a case split on an index) is refused with a reason; that work
belongs to \code{:prove}.

\subsection{Synthesis inside a proof}\label{sec:synthprove}

Bare \code{:synth} in prove mode targets the current goal \emph{with
its hypotheses as premises}. Candidates are offered as \code{it1},
\code{it2}, \dots\ exactly as in the session, except that here
\code{itN} splices the candidate applied to your hypotheses --- so
\code{exact it1} closes the goal:

\begin{session}
λ> :prove ∀ p q : Prop, (p → q) → p → q ∧ p
⊢ ∀ (p q : Prop), (p → q) → p → q ∧ p
⊢> intro p q h hp
p q : Prop
h : p → q
hp : p
⊢ q ∧ p
⊢> :synth
(synthesizing with hypotheses p q h hp as premises)
  it1  fun _ _ f x => ⟨f x, x⟩
⊢> exact it1
All goals accomplished
⊢> :qed mp_and
saved: theorem mp_and : ∀ p q : Prop, (p → q) → p → q ∧ p
\end{session}

Unlike \code{exact?}, which finds an \emph{existing} lemma, this
\emph{composes a new term} from the goal's own material --- a
constructive complement to the finisher tactics, needing no premise
database and no imports.

\subsection{Engines, budgets, and honesty about both}\label{sec:engines}

The default engine is Djinn: complete, terminating, and the only
source of refutations. \code{:set synth-engine exference} switches to
the ranked heuristic search, and \code{both} runs the two together ---
Djinn's candidates first, Exference's new ones appended, negative
verdicts only from the engine entitled to them. Exference runs under
explicit budgets (\code{:set synth-steps}, default 4096), and
truncation is always reported:

\begin{session}
λ> :set synth-engine both
synth engine: both
λ> :synth ((A → B → C) → (A → B) → A → C)
  it1  fun f g x => f x (g x)
note: exference: search truncated: step limit reached, queue limit pruned 31990
\end{session}

The pure searches answer in microseconds to milliseconds; the cost
center is backend verification, a few hundred milliseconds per
candidate. A wall-clock guard (default 20\,s,
\code{LEANT\_SYNTH\_TIMEOUT}) covers the quantified goals whose
bounded instantiation can widen the space; hitting it is reported as
``no answer,'' never as a verdict.

\subsection{What \code{:synth} will not do}\label{sec:limits}

Present limits, stated plainly: it will not synthesize recursion or
induction (recursive types contribute constructors only); it will not
analyze dependent structure (atoms are transported, never opened);
quantifier handling beyond the propositional fragment follows Djex's
bounded rank-N and guarded impredicativity rules --- sound and
terminating over an undecidable space, and widening as Djex improves
--- so some higher-rank goals answer ``no term found within bounds''
rather than anything stronger; and although Djex itself resolves
type-class evidence, Leant currently leaves instance arguments to
Lean's elaborator during verification, which is better placed to
resolve them. The fragment is bounded on purpose: inside it, answers
are fast, verified, and --- for refutations --- final.

%==============================================================================
\section{How it works, briefly}\label{sec:how}

A \code{:synth} query passes through five checked stages. A Lean
metaprogram (compiled once into a cached side environment) elaborates
the goal and serializes it into the engines' fragment: structural
connectives are decomposed; qualifying inductives expand into
constructor lists; everything else becomes an $\alpha$-normalized
opaque atom. The fragment translator either accepts the result or
refuses with a reason. The selected engine searches; its candidates
are rendered back into Lean syntax --- constructor names restored,
binders named by role, anonymous-constructor and leading-dot
spellings preferred with fully qualified fallbacks. Finally the
backend re-elaborates each candidate against the original goal, and
only survivors are shown and bound. The engine boundary is narrow by
design: the searches behind it are swappable, and nothing they emit is
believed until Lean has type-checked it.

%==============================================================================
\section{Status}\label{sec:status}

Leant is a working prototype under active development inside the
\emph{ProveIt} repository. The Haskell implementation
(\code{Tools/Leant}) is primary; the Python original
(\code{Tools/LeantPy}) remains as a lighter alternative without
synthesis. The synthesis design and its phased history are documented
in \code{Tools/Leant/SYNTHESIS\_PROPOSAL.md}; transcript-level golden
tests live in \code{Tools/Leant/test}. Expect sharp edges: the project exists to find out how far a
conversational, verification-gated synthesis workflow can be pushed,
and it moves at the speed of that question.
\end{document}
